\section{Introducción al estado del arte}

La detección de amenazas mediante aprendizaje automático cuenta con más de una década de desarrollo académico consolidado, particularmente en el dominio de la detección de intrusiones a nivel de red. Sin embargo, la aplicación de estas técnicas a la capa de aplicación de servicios web modernos --y en particular a arquitecturas asíncronas orientadas a eventos como Node.js-- constituye un área considerablemente menos explorada en la literatura indexada. Este capítulo revisa crítica y comparativamente los antecedentes académicos y técnicos directamente relacionados con el problema abordado por logSguarDian, organizados en cinco ejes: (1) los estudios de referencia sobre detección de intrusiones basada en aprendizaje automático, (2) los algoritmos base del modelo híbrido propuesto, (3) las metodologías de construcción de datasets etiquetados, (4) las restricciones de inferencia embebida documentadas en la literatura de Edge AI, y (5) las soluciones comerciales y perimetrales existentes junto con su respaldo documental. El objetivo no es enumerar cronológicamente estos trabajos, sino identificar, mediante su comparación directa, los vacíos metodológicos que este proyecto atiende.

\section{Detección de intrusiones basada en aprendizaje automático: panorama general}

Tres revisiones sistemáticas de amplio alcance delimitan el estado del arte general del uso de aprendizaje automático en ciberseguridad. \textcite{buczak2016} presentan una taxonomía comparativa de métodos de minería de datos y aprendizaje automático aplicados a la detección de intrusiones, evaluando su desempeño reportado en la literatura previa a 2016. Su hallazgo más relevante para este proyecto es doble: por un lado, confirman que los métodos de ensamble basados en árboles obtienen consistentemente una de las mejores relaciones entre precisión de clasificación y costo computacional de inferencia frente a otras familias de algoritmos; por otro, evidencian que la inmensa mayoría de los estudios revisados opera sobre datasets de flujo de red heredados (KDD Cup~99, NSL-KDD), centrados en características de nivel de transporte (duración de conexión, bytes transferidos, banderas TCP) y no en la semántica del contenido HTTP de la capa de aplicación.

\textcite{sarker2020} amplían este panorama proponiendo una taxonomía de "ciencia de datos de ciberseguridad" que organiza los enfoques de ML en función del tipo de amenaza abordada. Su revisión confirma la misma limitación estructural señalada por Buczak y Guven: los conjuntos de datos de referencia utilizados en la literatura mayoritaria siguen centrados en tráfico de red genérico, con escasa atención a la semántica de peticiones HTTP individuales o al comportamiento específico de APIs REST. \textcite{verma2020} conducen una revisión enfocada en sistemas de detección de intrusiones para entornos de Internet de las Cosas (IoT), comparando el desempeño de múltiples algoritmos --incluyendo Random Forest-- sobre datasets como NSL-KDD, CICIDS2017 y BoT-IoT. Su aporte principal es una tabla comparativa de exactitud por algoritmo y dataset que refuerza empíricamente la superioridad de los métodos de ensamble sobre alternativas más complejas en escenarios de recursos computacionales restringidos, un hallazgo directamente aplicable al contexto de inferencia embebida de logSguarDian, aunque su enfoque permanece anclado en dispositivos IoT y no en servidores de aplicaciones Node.js.

\textcite{berman2019} completan este panorama con una revisión específica de métodos de aprendizaje profundo aplicados a ciberseguridad, cubriendo arquitecturas como redes neuronales convolucionales y recurrentes en tareas de detección de malware, tráfico de red y correo no deseado. Su conclusión más relevante para este proyecto es una advertencia recurrente: los modelos de aprendizaje profundo de alta capacidad exhiben ganancias marginales de precisión frente a métodos más ligeros cuando se contabiliza el costo de inferencia y el tamaño del modelo serializado, lo cual constituye una justificación empírica adicional --proveniente de una fuente independiente a la arquitectura de este proyecto-- para descartar dichas arquitecturas en favor de modelos de ensamble ligeros, como se establece en la delimitación del alcance del componente de Machine Learning.

La comparación crítica de estos tres trabajos revela una convergencia clara: la literatura de referencia en detección de intrusiones mediante ML privilegia consistentemente los métodos de ensamble basados en árboles por su balance entre precisión y costo de inferencia, pero lo hace casi exclusivamente sobre datasets de flujo de red heredados que no capturan la semántica de una petición HTTP individual. Este vacío metodológico --la ausencia de literatura consolidada sobre feature engineering a nivel de request HTTP para APIs REST modernas-- es precisamente el espacio que ocupa el pipeline de extracción de características de logSguarDian, descrito en la sección 5.6 del protocolo de este trabajo.

\section{Algoritmos base del modelo híbrido: Random Forest e Isolation Forest}

El componente supervisado del modelo propuesto se fundamenta directamente en el trabajo seminal de \textcite{breiman2001}, quien introduce Random Forest como un método de ensamble que combina bagging con selección aleatoria de características en cada nodo de división, logrando robustez frente al sobreajuste sin sacrificar velocidad de inferencia. El componente no supervisado se fundamenta en \textcite{liu2008}, quienes proponen Isolation Forest bajo el principio de que las observaciones anómalas requieren, en promedio, menos particiones aleatorias para ser aisladas que las observaciones normales. Ambos algoritmos constituyen antecedentes directos --no meramente referenciales-- de la arquitectura del proyecto: sus propiedades de complejidad computacional lineal y ausencia de requisitos de normalización estricta son la razón técnica por la que fueron seleccionados frente a las quince alternativas evaluadas conceptualmente en la literatura revisada.

\textcite{chandola2009} contextualizan estos dos algoritmos dentro de una taxonomía más amplia de técnicas de detección de anomalías, estableciendo la distinción formal entre detección supervisada (que requiere ejemplos etiquetados de cada clase de anomalía) y detección no supervisada de una sola clase (que modela únicamente la normalidad). Su marco conceptual es el antecedente teórico directo de la decisión arquitectónica de combinar ambos paradigmas en un modelo híbrido, en lugar de optar por uno solo: mientras el primero garantiza precisión sobre las categorías de ataque conocidas, el segundo provee una capa complementaria de cobertura frente a variantes no vistas durante el entrenamiento.

El antecedente aplicado más reciente y más cercano a la arquitectura híbrida propuesta es el de \textcite{zeng2024}, quienes combinan técnicas de aprendizaje automático supervisado y no supervisado para la detección de anomalías en un contexto de "conciencia situacional de seguridad de red". Este trabajo confirma la viabilidad general de la combinación híbrida, pero opera nuevamente a nivel de flujo de red y no reporta restricciones de latencia ni de huella de memoria compatibles con un entorno de inferencia embebida como el de un middleware Node.js, dejando sin resolver la pregunta central de este proyecto: si dicha combinación híbrida es viable bajo las restricciones de tiempo real de un servidor de aplicaciones.

El antecedente de sistema más directamente comparable con el producto final de este proyecto es el de \textcite{chen2024}, quienes documentan un Web Application Firewall potenciado con inteligencia artificial orientado a proteger infraestructura crítica estadounidense contra exploits de día cero. Aunque comparte el objetivo general de mejorar la detección frente a las limitaciones de los WAF tradicionales basados en firmas, este trabajo mantiene la arquitectura perimetral clásica --el modelo opera como proxy inverso ante el servidor, no embebido en su proceso de ejecución-- y está diseñado para el contexto de infraestructura crítica de gran escala, sin reportar restricciones de recursos, latencia por petición ni disponibilidad de un artefacto de código abierto reproducible. Esta comparación evidencia que, incluso en el antecedente más cercano identificado, la propuesta de un modelo híbrido embebido directamente en el proceso de la aplicación --sin componente perimetral ni dependencia de infraestructura externa-- permanece sin precedente documentado.

\section{Metodologías de construcción y etiquetado de datasets}

La metodología de generación de tráfico sintético etiquetado adoptada para la construcción del corpus de entrenamiento de logSguarDian --descrita en la sección 5.10 del protocolo-- tiene como antecedente metodológico directo el trabajo de \textcite{shiravi2012}, quienes proponen un enfoque sistemático para generar datasets de referencia mediante el despliegue de un entorno de prueba controlado, la simulación de tráfico legítimo representativo y la ejecución de herramientas de ataque automatizadas, con etiquetado determinista basado en el conocimiento del proceso de generación ("ground truth by design"). Este enfoque, que dio origen a los datasets de la familia CICIDS ampliamente citados en la literatura de detección de intrusiones, constituye la base metodológica adoptada en este proyecto para garantizar la exactitud del etiquetado sin depender de anotación manual.

La comparación crítica con este antecedente revela, no obstante, la misma limitación estructural identificada en la sección anterior: los datasets resultantes de esta metodología --incluyendo el propio CICIDS-- se construyen predominantemente sobre características de flujo de red (duración, bytes por segundo, número de paquetes) y no sobre la semántica del contenido de la petición HTTP. La adopción del mismo principio metodológico de generación controlada, pero aplicado a la extracción de características semánticas del objeto \texttt{req} de Express en lugar de metadatos de flujo de red, representa una adaptación metodológica específica de este proyecto frente a su antecedente directo.

\section{Restricciones de inferencia embebida y Edge AI}

La ejecución de modelos de aprendizaje automático dentro de un proceso con restricciones estrictas de latencia y memoria --sin comunicación con servicios externos-- está documentada de forma consolidada en la literatura de Edge AI, cuyo texto de referencia es el de \textcite{warden2019}. Los autores establecen los principios de diseño para desplegar modelos ligeros bajo presupuestos de memoria y tiempo de inferencia del orden de milisegundos, incluyendo técnicas de cuantización y formatos de serialización optimizados como los adoptados en este proyecto (ONNX). Sin embargo, este cuerpo de literatura se desarrolla casi en su totalidad sobre microcontroladores y dispositivos embebidos de hardware dedicado, no sobre entornos de ejecución de un solo hilo orientados a eventos como el de Node.js. No se identificó en la revisión de literatura un antecedente académico que documente la aplicación de esta disciplina de inferencia embebida a las restricciones específicas del Event Loop de Node.js mediante \texttt{worker\_threads}, lo cual constituye una adaptación metodológica original de este proyecto sobre un cuerpo de conocimiento previamente establecido en un dominio de hardware distinto.

\section{Soluciones perimetrales existentes y su respaldo documental}

Como complemento a la revisión académica, la literatura técnica y los informes de la industria documentan de forma consistente las limitaciones estructurales de las soluciones perimetrales tradicionales frente a amenazas de capa de aplicación. La taxonomía del \textcite{owasp2021} establece las categorías de vulnerabilidad más críticas en aplicaciones web modernas --incluyendo inyección y fallos de control de acceso-- y sirve como marco de referencia frente al cual se evalúa la cobertura de logSguarDian. El reporte anual de \textcite{verizon2024} y el informe de \textcite{snyk2024} aportan evidencia cuantitativa reciente sobre la prevalencia de vulnerabilidades en cadenas de suministro de software abierto y sobre la efectividad limitada de los controles perimetrales frente a ataques dirigidos a la lógica de negocio. A diferencia de los antecedentes académicos revisados en las secciones anteriores, estas fuentes constituyen documentación técnica de la industria y no literatura indexada por pares; se incorporan en este capítulo únicamente como evidencia contextual de la magnitud del problema, no como fundamento metodológico del diseño del sistema, distinción que se mantiene consistente con la naturaleza aplicada de este trabajo.

\section{Síntesis comparativa}

La Tabla \ref{cuadro:sintesis_antecedentes} resume los antecedentes revisados en función de cuatro dimensiones críticas para el problema de investigación: el nivel de la pila de red en el que opera cada antecedente, la naturaleza del dataset o evidencia empleada, la arquitectura de despliegue del enfoque, y la limitación principal identificada frente a la propuesta de logSguarDian.

\begin{table}[h]
\small
\begin{tabular}{|p{2.7cm}|p{2.1cm}|p{2.9cm}|p{5.3cm}|}
\hline
\textbf{Antecedente} & \textbf{Nivel} & \textbf{Arquitectura} & \textbf{Limitación frente a logSguarDian} \\ \hline
Buczak \& Guven (2016); Sarker et al. (2020); Verma \& Ranga (2020) & Red (flujo) & Offline / batch & Datasets de flujo de red (KDD/NSL-KDD/CICIDS); sin semántica de petición HTTP \\ \hline
Berman et al. (2019) & Red / malware & Offline / batch & Modelos de aprendizaje profundo de alto costo, sin evaluación de restricciones de inferencia en tiempo real \\ \hline
Zeng et al. (2024) & Red (flujo) & Centralizada & Modelo híbrido supervisado/no supervisado sin restricciones de latencia embebida documentadas \\ \hline
Chen et al. (2024) & Aplicación & Perimetral (WAF/proxy) & Arquitectura de borde de red, no embebida en el proceso de la aplicación; sin artefacto reproducible \\ \hline
Shiravi et al. (2012) & Red (flujo) & Entorno de prueba controlado & Metodología de generación adoptada, pero features de flujo, no de contenido HTTP \\ \hline
Warden \& Situnayake (2019) & Hardware embebido & Microcontroladores & Disciplina de Edge AI no documentada para el Event Loop de Node.js \\ \hline
\end{tabular}
\caption[Síntesis comparativa de antecedentes]{Síntesis comparativa de los antecedentes revisados frente a la propuesta de logSguarDian. Elaboración propia a partir de la revisión de literatura de este capítulo.}
\label{cuadro:sintesis_antecedentes}
\end{table}

\section{Conclusión: vacíos identificados y aporte del proyecto}

La revisión crítica de los antecedentes académicos y técnicos converge en tres vacíos concretos que este proyecto atiende de forma directa. En primer lugar, la literatura de referencia en detección de intrusiones mediante aprendizaje automático --representada por las revisiones de Buczak y Guven, Sarker et al. y Verma y Ranga-- privilegia consistentemente los datasets de flujo de red heredados sobre la semántica de peticiones HTTP individuales, dejando sin resolver la pregunta de cómo construir un pipeline de ingeniería de características orientado específicamente al objeto de petición de un framework como Express. En segundo lugar, el antecedente aplicado más cercano en arquitectura híbrida supervisada/no supervisada (Zeng et al., 2024) y el antecedente de sistema más cercano en propósito (Chen et al., 2024) coinciden en operar bajo arquitecturas perimetrales o centralizadas, sin abordar las restricciones específicas de un despliegue embebido dentro del proceso de una aplicación Node.js con recursos limitados, que es precisamente el escenario documentado para las startups tecnológicas guatemaltecas en el planteamiento del problema de este trabajo. En tercer lugar, la disciplina de inferencia embebida bajo restricciones de latencia y memoria, consolidada en la literatura de Edge AI para hardware de microcontroladores, no cuenta con un antecedente documentado de aplicación al modelo de concurrencia de un solo hilo de Node.js mediante \texttt{worker\_threads}.

Estos tres vacíos --ausencia de feature engineering semántico a nivel de HTTP en la literatura de ML-IDS, ausencia de arquitecturas híbridas embebidas (no perimetrales) documentadas para aplicaciones web, y ausencia de una adaptación de la disciplina de Edge AI al entorno de ejecución de Node.js-- no son independientes entre sí: se articulan en un mismo vacío estructural, la falta de una solución de detección de amenazas que combine aprendizaje automático híbrido con inferencia embebida en tiempo real dentro de un servidor Node.js, sin infraestructura perimetral ni dependencias externas. logSguarDian se posiciona como una respuesta directa a este vacío convergente, adoptando de sus antecedentes los elementos validados empíricamente por la literatura --la superioridad de los métodos de ensamble ligeros sobre alternativas de mayor capacidad, la complementariedad teórica entre detección supervisada y no supervisada, y la metodología de generación controlada de datasets con etiquetado determinista-- y extendiéndolos hacia un dominio de aplicación, la protección embebida de APIs Node.js, en el que dicha combinación no había sido previamente documentada.
